% =====================================================================
%  Выводы по лабораторной работе.
% =====================================================================

\labpart{Выводы}

В ходе выполнения лабораторной работы реализована система
автоматического машинного перевода текстов по варианту 11
(направление англо-французское, предметная область -- научные статьи
по медицине и критика предметов изобразительного искусства) как новый
сценарий использования поверх архитектуры, разработанной в ЛР1--ЛР3.
В отличие от ЛР2 и ЛР3, для перевода не понадобился отдельный
микросервис: пословный перевод по словарю и синтаксический разбор
предложения -- такие же операции без состояния над spaCy, как уже
существующие в \texttt{nlp-service} токенизация и лемматизация, поэтому
они естественно расширили этот сервис, а не потребовали нового; ни
один существующий контракт системы (поиск, краулинг, определение
языка, реферирование) при этом не был нарушен.

Реализована система прямого (пословного) перевода: каждое слово текста
ищется в двуязычном словаре по паре (лемма, часть речи) с откатом на
запись «любая часть речи», если точное совпадение не найдено -- откат
компенсирует то, что словарь размечен частями речи вне контекста
конкретного предложения. Определяющим фактором качества такой системы
является размер словаря, а не сам алгоритм подстановки: словарь
объединяет стартовый курируемый список (311 статей) с реальным
двуязычным словарём MUSE, размеченным частями речи через spaCy (итог
-- 87~791 статья, без ограничения по частоте исходного слова), и даёт
среднее покрытие 91{,}0\,\% на реальной тестовой коллекции
(табл.~\ref{tab:dictionary-coverage}). Отдельно решена
сопутствующая проблема корректности: значение «любая часть речи» в
столбце \texttt{pos} хранится явным значением-часовым, а не
\texttt{NULL} -- в PostgreSQL \texttt{NULL} не равен самому себе для
уникального ограничения, из-за чего утилита пополнения словаря могла
бы молча копить противоречивые дубликаты записей для одного и того же
слова.

По аналогии с тестовыми прогонами ЛР2 (\texttt{lang\_id\_runs}) и ЛР3
(\texttt{summarization\_runs}) добавлена вкладка «Тестирование»:
перевод коллекции одним нажатием, живой прогресс-бар через WebSocket
и средние метрики (число слов до и после перевода, покрытие словаря,
время) на странице «Метрики». Тестовый прогон использует уже
существующую в системе разметку языка документа
(\texttt{documents.language}, определённую при кроулинге -- см.
ЛР1/ЛР2) и переводит только документы на исходном языке перевода;
документы на других языках коллекции в прогон не попадают, поскольку
их слова почти никогда не находятся в словаре по построению (поиск
идёт по лемме «чужого» языка) и только искажали бы среднее покрытие
словаря.

Обе дополнительные вкладки задания используют уже существующую в
проекте функциональность: частотный список слов (вкладка~1)
переиспользует тот же конвейер токенизации с исключением стоп-слов,
что и индексация документов ЛР1, а дерево синтаксического разбора
(вкладка~2) -- перенесённый и адаптированный из проекта прошлого
семестра (\texttt{corpus-manager}) конвейер spaCy с включённым
зависимостным парсером (в основном конвейере индексации он отключён
ради скорости). POS-теги и синтаксические роли выводятся как сырые
значения spaCy без перевода на русский язык -- решение, взятое из
исходного \texttt{corpus-manager} и сознательно сохранённое, поскольку
расшифровка тега поверх уже переведённого текста избыточна для
читателя, знакомого с разметкой Universal POS. Отдельная утилита
пополнения/корректировки словаря (постраничная таблица с поиском,
добавлением, редактированием и удалением статей) была фактически
задействована в ходе самой лабораторной работы -- отсутствующие в
словаре слова предметной области (например, специализированные
медицинские термины) добавлялись в БД без переразвёртывания системы.
